%%
%% Arquivo principal do TCC.
%%
%% Siga as instruções que seguem comentadas neste arquivo.
%%
%% ESTA VERSÃO REQUER UM EDITOR UTF-8.
%%
%% Exemplos de utilização dos comandos podem ser encontrados
%% nos capítulos que acompanham este pacote.
%%
%% 
\documentclass[tcc]{tcc}

% algumas customizações: o preamble.tex pode ser editado pelo usuário
\input{preamble}

% modifique os argumentos dos comandos abaixo para que as primeiras páginas
% sejam construídas automaticamente
% \title{Construção de um banco de dados de perfis sigma para modelos tipo COSMO}
\title{Título do trabalho}
\author{Nome Autor}
\authorCitation{Autor, Nome.}
\tyear{2019} 
% verificar o número de páginas do PDF final para inclusão da informação na ficha
% catalográfica!
\totalPages{160}
\advisor{Nome Orientador Um, D.Sc.} 
\advisorCitation{Um, Nome O.}
%%%%%%% SE TIVER 2 ORIENTADORES, DESCOMENTAR!
% \advisori{Nome Orientador Dois, D.Sc.}
% \advisoriCitation{Dois, Nome O.} 
%%%%% COMENTAR SE NÃO TIVER COORIENTADOR
\coadvisor{Nome Coorientador, D.Sc.}
\coadvisorCitation{Coorientador, Nome.}


\agradecimento{
Agradeço ........}

\resumo{ Incluir aqui o texto do resumo}
\palavraschave{palavra chave 1, palavra chave 2, palavra chave 3, palavra chave 4, palavra chave 5}

% Desse ponto modifique apenas a inclusão dos capítulos e apêndices
\begin{document}
% hifenização
\include{hyphen}

\maketitle

% no arquivo CpaSymbols.tex devem ser adicionados os símbolos citados
% nas equações do trabalho
\include{CapSymbols}

% Inclusão dos Capítulos conforme a estrutura definida/escolhida para o trabalho!
\include{Cap1}
\include{Cap2}
\include{Cap3}  
\include{Cap4}
\include{Cap5} 

%%%%% NÃO ALTERAR!  
% Comandos para as referencias bibliográficas
\newdimen\bibindent
\setlength\bibindent{1.5em} % identaçao das referencias
% Um lineskip menor neste contexto de referencias
\baselineskip 4.3mm
% Edite o arquivo bib.bib com as suas próprias referências 
\bibliography{bib}


% Inclusao dos Appendices, se for necessário! Criar o/os arquivos .tex 
 \appendix
%\include{Ap1}
% \include{Ap2}
% \include{Ap3}

\end{document}
